\documentclass[12pt,a4paper,fontset=none]{ctexart}

\usepackage[a4paper,margin=2.5cm]{geometry}
\usepackage{enumitem}
\usepackage{hyperref}
\usepackage{longtable}
\usepackage{xcolor}

\setCJKmainfont[
  Path=/System/Library/Fonts/,
  BoldFont={STHeiti Medium.ttc}
]{STHeiti Light.ttc}
\setCJKsansfont[
  Path=/System/Library/Fonts/,
  BoldFont={STHeiti Medium.ttc}
]{STHeiti Light.ttc}
\setCJKmonofont[
  Path=/System/Library/Fonts/
]{STHeiti Light.ttc}

\hypersetup{
  colorlinks=true,
  linkcolor=blue,
  urlcolor=blue
}

\setlist[itemize]{leftmargin=2em}
\setlist[enumerate]{leftmargin=2em}
\renewcommand{\arraystretch}{1.25}

\title{LLM-Enhanced Multimodal Product Recommendation\\
with Review Aspect-Specific Knowledge Graphs}
\author{RA-GARK 專案交接文件}
\date{\today}

\begin{document}
\maketitle

\section{Environment}

\begin{itemize}
  \item 建議環境：\texttt{py311\_cuda126}。
  \item 所有 Python 指令皆從 \texttt{Code/} 目錄執行。
  \item 目前 repo 沒有固定提供完整 environment file；若在新機器接手，請先安裝 Python 依賴，再執行 smoke test 確認環境與資料路徑。
\end{itemize}

\subsection*{程式碼架構}

\begin{verbatim}
Code/
    config.py
    data.py
    kg_loader.py
    model.py
    losses.py
    evaluate.py
    utils.py
    train_ragark.py
    run_ablations.py
    run_main_benchmark.py
    case_study.py
    smoke_test.py
    baselines/
    data/
    handoff/
    relation_init/
    results/
    runs/
\end{verbatim}

\begin{longtable}{p{0.30\textwidth}p{0.60\textwidth}}
\textbf{檔案或資料夾} & \textbf{說明} \\
\hline
\texttt{config.py} & 實驗設定檔，包含 data path、模型參數、訓練參數與 ablation flags。 \\
\texttt{data.py} & 載入 review interaction、建立 KG index、建立 LightGCN adjacency。 \\
\texttt{kg\_loader.py} & canonical KG loader；只有 canonical KG branch 或檢查工具會用到。 \\
\texttt{model.py} & RA-GARK 模型架構。 \\
\texttt{losses.py} & BPR loss、aspect-level contrastive loss、user cross-view contrastive loss。 \\
\texttt{evaluate.py} & Top-K recommendation evaluation。 \\
\texttt{train\_ragark.py} & RA-GARK 主模型訓練與測試。 \\
\texttt{run\_ablations.py} & 消融實驗與 sensitivity 實驗。 \\
\texttt{run\_main\_benchmark.py} & baseline 與 RA-GARK 主比較實驗。 \\
\texttt{case\_study.py} & rationale attention case study。 \\
\texttt{smoke\_test.py} & 快速檢查 data、KG、model forward/backward、evaluation 是否能跑通。 \\
\texttt{baselines/} & KGAT、KGCL、KGRec、MCCLK baseline。 \\
\texttt{relation\_init/} & canonical KG initialization 的 archived experiment。 \\
\texttt{results/} & 最新輸出結果，例如 main benchmark、ablation、case study。 \\
\texttt{runs/} & checkpoint、best ledger、archive CSV。 \\
\end{longtable}

\subsection*{Smoke Test}

使用方式：

\begin{verbatim}
cd Code
python smoke_test.py
\end{verbatim}

\texttt{smoke\_test.py} 的目的只是確認環境與主要 code path 可正常執行。它不是正式實驗，也不是論文數字來源。

\section{Datasets}

\begin{itemize}
  \item 主要資料來源：Amazon Reviews 2023。
  \item 使用類別：Books。
  \item 原始 review 來源：\url{https://amazon-reviews-2023.github.io/}
  \item 目前主實驗使用的是已 preprocess 完成的檔案，位於 \texttt{Code/data/}。
\end{itemize}

\subsection*{目前實際使用的 data}

\begin{longtable}{p{0.38\textwidth}p{0.52\textwidth}}
\textbf{檔案} & \textbf{用途} \\
\hline
\texttt{data/reviews\_30\_20.pkl} & 主要 interaction data。所有 RA-GARK、baseline、case study 預設都會讀這個檔案。 \\
\texttt{data/df\_edges\_item\_aspect1.csv} & 預設 item-aspect KG，對應 \texttt{Config.kg\_path}。 \\
\texttt{data/df\_edges\_all\_aspect1.csv} & raw KG input，供 \texttt{kg\_clean.py} 產生 canonical KG。 \\
\texttt{data/kg\_canonical.csv} & canonical KG；只有 \texttt{use\_canonical\_kg=True} 或 \texttt{relation\_init/} 會使用。 \\
\texttt{data/kg\_relation\_map.json} & raw relation 到 canonical relation 的 mapping，用於 audit。 \\
\texttt{data/kg\_canonical\_report.txt} & canonical KG 清理報告。 \\
\end{longtable}

\subsection*{主要 config 參數}

\begin{longtable}{p{0.30\textwidth}p{0.42\textwidth}p{0.18\textwidth}}
\textbf{參數名稱} & \textbf{說明} & \textbf{值範例} \\
\hline
\texttt{interaction\_path} & preprocess 後的 review interaction 檔案。 & \texttt{data/reviews\_30\_20.pkl} \\
\texttt{kg\_path} & 預設 item-aspect KG。 & \texttt{data/df\_edges\_item\_aspect1.csv} \\
\texttt{canonical\_kg\_path} & canonical KG 檔案。 & \texttt{data/kg\_canonical.csv} \\
\texttt{use\_canonical\_kg} & 是否改用 canonical KG。預設為 \texttt{False}。 & \texttt{False} \\
\texttt{kg\_top\_freq\_pct} & legacy KG 中要剪掉的高頻 aspect 比例。 & \texttt{0.02} \\
\texttt{num\_aspects} & RA-GARK aspect slots 數量。 & \texttt{4} \\
\texttt{embedding\_dim} & embedding 維度。 & \texttt{128} \\
\texttt{eval\_k} & evaluation 的主要 Top-K。 & \texttt{20} \\
\end{longtable}

\section{Preprocess}

目前 repo 內保留的是已 preprocess 完成、可直接訓練的資料。正式接手時請先確認 \texttt{Code/data/} 內上述資料皆已準備完成，再依照以下順序執行。

\subsection*{Preprocess 做的事情}

目前主訓練資料 \texttt{reviews\_30\_20.pkl} 已完成以下處理：

\begin{itemize}
  \item 過濾 review interaction，只保留可用欄位。
  \item 依照 like 欄位保留 positive interactions 作為推薦訓練資料。
  \item 對 user 與 item 重新 label encoding，產生 \texttt{user\_idx} 與 \texttt{item\_idx}。
  \item 使用 user-stratified split 產生 train、validation、test。
  \item 訓練時使用 \texttt{data/df\_edges\_item\_aspect1.csv} 建立 item-aspect KG index。
  \item KG 會剪掉高頻 noisy aspects 與 \texttt{Config.kg\_stopwords} 中的 stopwords。
\end{itemize}

如果要重新整理 canonical KG，可執行：

\begin{verbatim}
python kg_clean.py
\end{verbatim}

此步驟會以 \texttt{data/df\_edges\_all\_aspect1.csv} 為 input，輸出：

\begin{itemize}
  \item \texttt{data/kg\_canonical.csv}
  \item \texttt{data/kg\_relation\_map.json}
  \item \texttt{data/kg\_canonical\_report.txt}
\end{itemize}

\subsection*{1. 執行 RA-GARK}

執行主模型：

\begin{verbatim}
python train_ragark.py
\end{verbatim}

此指令會使用 \texttt{config.py} 中的預設設定進行訓練與 test evaluation。預設流程包含：

\begin{itemize}
  \item 載入 \texttt{reviews\_30\_20.pkl}
  \item 建立 item-aspect KG index
  \item 建立 LightGCN adjacency
  \item 使用 RA-GARK 訓練推薦模型
  \item 以 validation NDCG 選擇 best checkpoint
  \item 載入 best checkpoint 並輸出 test metrics
\end{itemize}

\subsection*{2. 執行消融實驗}

快速檢查版：

\begin{verbatim}
python run_ablations.py --mode minimal --reuse
\end{verbatim}

論文表格版本：

\begin{verbatim}
python run_ablations.py --mode paper --reuse
\end{verbatim}

sensitivity 實驗：

\begin{verbatim}
python run_ablations.py --mode sensitivity --reuse
\end{verbatim}

主要輸出：

\begin{itemize}
  \item \texttt{results/ablation\_results\_<mode>.csv}
  \item \texttt{runs/archive/ablation\_results\_<mode>\_<timestamp>.csv}
  \item \texttt{runs/best/}
  \item \texttt{runs/best\_ledger.json}
\end{itemize}

\subsection*{3. 執行 Baseline 模型}

各 baseline 由 \texttt{run\_main\_benchmark.py} 統一執行：

\begin{verbatim}
python run_main_benchmark.py
\end{verbatim}

包含模型：

\begin{itemize}
  \item MCCLK
  \item KGCL
  \item KGAT
  \item KGRec
  \item LightGCN
  \item RA-GARK
\end{itemize}

主要輸出：

\begin{itemize}
  \item \texttt{results/main\_benchmark\_results.csv}
  \item \texttt{runs/archive/main\_benchmark\_results\_<timestamp>.csv}
\end{itemize}

\subsection*{4. Case Study}

執行：

\begin{verbatim}
python case_study.py
\end{verbatim}

用途：輸出 RA-GARK softmax rationale attention 的案例，觀察同一商品在不同 user 下的 aspect attention 是否不同。

主要輸出：

\begin{itemize}
  \item \texttt{results/case\_study.txt}
  \item \texttt{results/case\_study.csv}
\end{itemize}

\subsection*{5. 結果檔案對應}

\begin{longtable}{p{0.38\textwidth}p{0.52\textwidth}}
\textbf{檔案} & \textbf{說明} \\
\hline
\texttt{results/main\_benchmark\_results.csv} & 主模型與 baseline 的比較結果。 \\
\texttt{results/ablation\_results\_paper.csv} & 論文消融實驗表格主要來源。 \\
\texttt{results/ablation\_results\_full.csv} & 較完整的消融結果。 \\
\texttt{results/case\_study.txt} & human-readable case study。 \\
\texttt{results/case\_study.csv} & machine-readable case study attention weights。 \\
\texttt{runs/best\_ledger.json} & 各 config 的 best-ever metrics ledger。 \\
\end{longtable}

\subsection*{注意事項}

\begin{itemize}
  \item 不要任意覆蓋 \texttt{runs/best/} 或 \texttt{runs/best\_ledger.json}，除非明確要重新更新 best-run ledger。
  \item \texttt{relation\_init/} 是 archived experiment，除非要重訪 canonical KG initialization branch，否則不建議修改。
  \item 若只想確認環境，請執行 \texttt{smoke\_test.py}；若要取得正式實驗數字，請執行 \texttt{run\_ablations.py} 或 \texttt{run\_main\_benchmark.py}。
  \item 編譯本交接文件：
\end{itemize}

\begin{verbatim}
cd Code/handoff
latexmk -xelatex -interaction=nonstopmode main.tex
\end{verbatim}

\end{document}
